\subsection{Claim-level provenance and verification systems}

There is also direct prior art for representing and verifying claims at finer granularity than whole documents. Zhang, Ives and Roth formalize provenance graphs for natural-language claims in order to track where a claim came from and how it evolved \cite{zhang2020claim}. In biomedicine, the SEE framework represents scientific claims together with their provenance, argumentative relations, subjects, methods and assumptions in an RDF/OWL evidence model \cite{boelling2014see}. These systems make clear that claim-level provenance and evidence structure are not new functions introduced by RAKL.

Verification systems add another layer. ProVe checks whether knowledge-graph triples are actually supported by the textual sources recorded as their provenance rather than assuming that a documented source entails the graph assertion \cite{amaral2024prove}. Open-domain scientific claim verification further shows that retrieval source and retrieval method materially affect evidence recall and precision when the relevant document is not already supplied to the verifier \cite{vladika2024}. In scholarly LLM interfaces, PaperTrail decomposes generated answers and sources into claim/evidence units to surface unsupported assertions and omissions \cite{martinboyle2026}.

End-to-end research systems now make the claim--evidence boundary still less available as a novelty claim. ScientistOne maintains a chain of evidence across literature review, solution discovery and paper writing, and couples it to post-hoc checks for score verification, specification violations, references and method--code alignment \cite{meng2026scientistone}. A 2026 survey of autonomous research agents likewise identifies claim verification and auditability, rather than mere workflow execution or code release, as a central unresolved systems gap \cite{ding2026verificationgap}. These works own substantial territory around evidence-traceable research generation and reviewer-facing audit. The residual object here is therefore not a chain-of-evidence mechanism by itself, but the typed scientific-state semantics governing what such evidence is allowed to change, under which contexts and authority coordinates, and how incompatible or later-revoked states remain represented.

Research-agent trajectories now supply an even closer operational parent. Zhuang et al. treat a completed experimental trajectory as candidate evidence rather than as knowledge, verify consequential artifacts before release, qualify cross-round intervention claims by execution validity and attribution, and retain only bounded Repairs or Guards while withholding unsupported findings \cite{zhuang2026trajectory}. Their industrial study also reports non-monotone adaptation---for 22 of 30 methods the last valid round was worse than an earlier best---and identifies affirmative applicability judgment as a remaining controller bottleneck. This owns substantial territory around verified proposal handoffs, negative-result retention, applicability-bounded experimental records and non-monotone research history. It also provides empirical motivation for preserving validated intermediate states rather than treating the final trajectory endpoint as authoritative.

A complementary 2026 evaluation tests the epistemic process itself rather than only task outcomes. R\'ios-Garc\'ia et al. evaluate scientific agents across multiple domains and report systematic failures to use contradictory evidence and to revise beliefs under refutation even when outcome metrics appear successful \cite{riosgarcia2026}. Their result directly sharpens the design target here: successful workflow execution is not a certificate that the process respected evidence, and therefore outcome success cannot be allowed to write scientific authority by itself. Epistemic mechanics does not claim to repair the model-level reasoning deficits diagnosed in that study; it supplies a fail-closed state-governance layer intended to prevent those deficits from silently becoming canonical scientific state.

These works narrow the present contribution in a useful way. Epistemic mechanics does not claim novelty for claim decomposition, provenance graphs, evidence attribution, source retrieval, support/refute classification, chain-of-evidence auditing, bounded artifact verification or applicability-qualified experimental records in isolation. Its additional object is the broader governed scientific-state transition around those operations: what kind of support was observed, under what context and independence assumptions, which scientific authority coordinates it can change, whether higher-order alignment obstructions and conflicting objects remain explicitly represented, and whether later discovery can reopen a state that had previously been judged saturated. Conversely, the present formal paper does not reproduce the empirical validation of these research-agent systems or claim comparable empirical superiority.
